\section{Introduction}\label{sec:intro}

A securities regulator that wants to find manipulation faces three problems at once. It rarely has labelled cases, because enforcement is slow and selective. It cannot easily evaluate a rule change before introducing it, because the behaviour of manipulators changes with the rule. And the data that would let an outside researcher test a method, orders tagged with an account, is held only by the exchange and the regulator. These problems are acute in Vietnam. In ten recent decisions of the State Securities Commission with exact dates (Appendix~\ref{app:cases}) the manipulation was carried out with between 5 and 164 accounts (median 24.5), by trading continuously, or crossing orders among accounts, to create the appearance of supply and demand. Spoofing and layering, the focus of most of the detection literature \citep{tuccella2021,fabre2025}, do not appear among these cases.

This paper asks two questions of a regulator in that position and answers them with a combination of theory, simulation and public data. \emph{Can a ring of accounts be found without labelled cases?} And \emph{which rules and how much inspection capacity make a ring that adapts to them give up?} For the first we study a label-free coincidence test: each account's submissions are compared with the aggregate of all other accounts, and the comparison is calibrated by circularly shifting the account's own series, which gives an exact $p$-value under a stated exchangeability condition. For the second we let a ring re-optimise against each rule in a family of simulated markets that match published and measured moments of the Ho Chi Minh Stock Exchange, and we relate the answer to closed-form conditions.

We make six contributions. (i)~We state the test, prove that it controls false alarms exactly per account (Proposition~\ref{prop:exact}), and show that ranking many accounts at once cannot inherit that guarantee because of the resolution of the shift distribution (Proposition~\ref{prop:budget}). (ii)~We derive a law for its power, in which the shift of the statistic falls like $N^{-1/2}$ in the number of accounts and rises with the number of co-active members (Proposition~\ref{prop:power}), verify it by Monte Carlo and show that it explains the dependence on the number of accounts in the agent-based market; we show why netting buys against sells protects against honest market-making desks (Proposition~\ref{prop:netting}). (iii)~We turn the law into a deterrence condition for a regulator with a fixed inspection capacity and a fine proportional to the gain (Proposition~\ref{prop:deter}), and derive and test a law for a rule that prices off-book blocks at a trailing mean (Proposition~\ref{prop:block}). (iv)~We evaluate rules by the best response of a ring over two families of six calibrated markets, with the effect of rules on honest participants (Section~\ref{sec:policy}). (v)~We stress-test the statistic on real account-level order flow from a public venue, where it fails the exchangeability condition for a large share of accounts and detects a planted ring less often than in the simulator (Section~\ref{sec:real-accounts}), and we check what public data on Vietnam can say: documented enforcement cases, a natural experiment on the tick size, and pump events (Section~\ref{sec:real}). (vi)~We release the code and all result tables.

The paper does not claim that the detector works on real Vietnamese order data. Its manipulators and features are written by us, a supervised model trained in the simulator does better than the label-free test in every simulated scenario, the real account-level data come from a different kind of market, and the Vietnamese checks are at the level of the stock, not of the account. The contribution is a calibrated screening tool with a measured failure profile, a theory that links its power to deterrence capacity, and an evaluation method for rules that does not assume manipulators stand still. Section~\ref{sec:related} places the paper among the related work, Section~\ref{sec:sim} sets up the market and the regulator's problem, Section~\ref{sec:method} gives the method and its theory, Sections~\ref{sec:detect}--\ref{sec:real} report the experiments, and Section~\ref{sec:discussion} discusses the impact for policy and the limits.
